\documentclass[11pt]{article}
\usepackage[margin=1in]{geometry}
\usepackage{amsmath,amssymb,booktabs,siunitx,hyperref,microtype,xcolor}
\usepackage{enumitem}
\usepackage{longtable}
\usepackage{array}
\usepackage{listings}
\usepackage{fancyvrb}
\setlist[itemize]{leftmargin=*,topsep=2pt,itemsep=2pt}
\setlist[enumerate]{leftmargin=*,topsep=2pt,itemsep=2pt}

\hypersetup{
  colorlinks=true,
  linkcolor=black,
  citecolor=black,
  urlcolor=blue!60!black,
  pdftitle={Literature Provenance: PNP-BV ORR Forward Solver},
  pdfauthor={Jake Weinstein}
}

\lstdefinestyle{bibstyle}{
  basicstyle=\ttfamily\footnotesize,
  breaklines=true,
  breakatwhitespace=true,
  columns=fullflexible,
  frame=single,
  framesep=4pt,
  xleftmargin=4pt,
  xrightmargin=4pt
}

\newcommand{\true}{\mathrm{TRUE}}
\newcommand{\Eeq}{E_{\mathrm{eq}}}
\newcommand{\VRHE}{V_{\mathrm{RHE}}}
\newcommand{\VT}{V_{T}}
\newcommand{\Hpe}{\mathrm{H}^{+}}
\newcommand{\HtwoOtwo}{\mathrm{H}_2\mathrm{O}_2}
\newcommand{\ClOf}{\mathrm{ClO}_4^{-}}

\title{Literature Provenance: PNP--BV ORR Forward Solver\\
       Three Numerical Changes, Cited Precedents, and Genuine-Novelty Assessment}
\author{Jake Weinstein}
\date{Week of April~27, 2026}

\begin{document}
\maketitle

\begin{abstract}
This document audits the bibliographic foundations of the PNP--BV ORR forward
solver rebuild summarized in \texttt{PNP Inverse Solver Revised.tex}. For each
of the three stacked numerical changes ((1)~3-species PNP $+$ analytic
Boltzmann counterion, (2)~log-concentration $u_i=\ln c_i$ primary unknown,
(3)~log-rate Butler--Volmer evaluation) we identify the canonical literature
precedents, distinguish reuse from genuine novelty, flag and correct three
factually wrong citations in the existing bibliography (one of which is
fabricated), and recommend additions. Surrounding methods (ORR mechanism,
Tafel-ridge identifiability, charge/voltage continuation, adjoint sensitivity)
are also covered for reader background. All citations have been verified at
DOI/abstract level except those explicitly marked.
\end{abstract}

%% ============================================================
\section{Executive Summary}
%% ============================================================

\subsection*{Verdict table}

\begin{center}\small
\renewcommand{\arraystretch}{1.4}
\begin{tabular}{@{}p{0.16\textwidth} p{0.30\textwidth} p{0.27\textwidth} p{0.21\textwidth}@{}}
\toprule
\textbf{Change} & \textbf{Canonical precedent} & \textbf{Where the implementation deviates} & \textbf{Novelty assessment} \\
\midrule
\#1 PBNP hybrid &
Zheng \& Wei, \emph{J.~Chem.~Phys.}\ \textbf{134}, 194101 (2011),
\href{https://doi.org/10.1063/1.3581031}{doi:10.1063/1.3581031}; deeper
Gouy--Chapman--Stern (1910/1913/1924) and the Eisenberg--Coalson--Kurnikova--Lu
ion-channel PNP lineage. &
Transplanted into a PEMFC-style ORR solver in \SI{0.1}{\mole\per\liter}
$\mathrm{HClO}_4$ with spectator $\ClOf$ as the analytic-Boltzmann species.
No surveyed ORR/RDE/fuel-cell paper does this. &
\textbf{Modest.} Construction reused; \emph{application} to ORR with explicit
Boltzmann-$\ClOf$ inside a PNP--BV solver appears new in the surveyed
electrochemistry literature. \\
\midrule
\#2 log-density $u=\ln c$ &
Metti, Xu, Liu, \emph{J.~Comput.~Phys.}\ \textbf{306}, 1--18 (2016),
\href{https://doi.org/10.1016/j.jcp.2015.10.053}{doi:10.1016/j.jcp.2015.10.053};
high-order extension Fu \& Xu, \emph{CMAME} \textbf{395}, 115031 (2022),
\href{https://doi.org/10.1016/j.cma.2022.115031}{doi:10.1016/j.cma.2022.115031};
semiconductor ancestry through Slotboom 1969 and Brezzi--Marini--Pietra 1989. &
Pushes log-density into a \emph{full PNP--BV ORR solver}, with the same CG
order on $u_i$ and $\varphi$, monolithic Newton, backward Euler --- and
propagates $u_i$ through the BV boundary condition (Change~\#3), not just bulk
transport. &
\textbf{Genuine domain transfer.} No peer-reviewed ORR/RDE/fuel-cell PNP--BV
solver in the surveyed literature uses $\ln c$ as a primary Galerkin
unknown. \\
\midrule
\#3 log-rate BV &
Tafel 1905, \emph{Z.~Phys.~Chem.}\ \textbf{50U}, 641--712,
\href{https://doi.org/10.1515/zpch-1905-5043}{doi:10.1515/zpch-1905-5043}
(high-$\eta$ physical \emph{approximation}, drops a branch);
Fattal \& Kupferman, \emph{J.~Non-Newtonian Fluid Mech.}\ \textbf{123},
281--285 (2004),
\href{https://doi.org/10.1016/j.jnnfm.2004.08.008}{doi:10.1016/j.jnnfm.2004.08.008}
(closest \emph{structural} cousin: log-conformation tensor at high Weissenberg
number); analogous CHEMKIN PLOG / log-sum-exp / softmax patterns. &
The exact construction
$r=\exp(\ln k_0+u_{\mathrm{cat}}+\sum_f\nu_f(u_{\mathrm{sp}(f)}-\ln c_{\mathrm{ref},f})-\alpha n_e\eta/\VT)$,
with the cathodic stoichiometric-power loop \emph{inside} the log, evaluated
then exponentiated once, does not appear in any electrochemistry paper
located. The log-rate path also bypasses the symmetric \texttt{\_U\_CLAMP}
underflow protection. &
\textbf{Small local numerical novelty (defensible).} Algebraically trivial;
conceptually parallel to Fattal--Kupferman 2004 and CHEMKIN PLOG;
\emph{not documented} as a numerical strategy in any PNP--BV / ORR / RDE
paper located. \\
\bottomrule
\end{tabular}
\end{center}

\subsection*{Three corrections the existing LaTeX bibliography must make}

\begin{enumerate}
\item \textbf{The ``Liu, Maxian, Shen, Zitelli, Waltz, Monteiro 2022, CMAME 396,
      115070, arXiv:2105.01163'' citation is fabricated.}\ \ The arXiv ID
      resolves to \textbf{Fu \& Xu, \emph{CMAME} \textbf{395}, 115031 (2022),
      doi:10.1016/j.cma.2022.115031} (two authors, not six; vol.~395 not 396;
      article 115031 not 115070). Replace.
\item \textbf{The Slotboom citation as printed is a chimera} --- wrong year for
      the title given, wrong title for the year given. The canonical
      Slotboom-variables reference is \textbf{Slotboom 1969, \emph{Electronics
      Letters} \textbf{5}(26), 677--678, doi:10.1049/el:19690510}. The
      ``PN-product'' paper, if cited, is Slotboom 1977, \emph{Solid-State
      Electron.}\ \textbf{20}(4), 279--283. Slotboom variables
      ($u=c\exp(-z\psi/\VT)$) and log-density ($u=\ln c$) are \emph{related}
      but \emph{distinct} substitutions; the writeup should not equate them.
\item \textbf{The ``Zheng-Chen-Wei JCP 230 (2011)'' cited as the PBNP source is
      the full-PNP second-order solver paper, not the PBNP paper.} The
      canonical PBNP paper is \textbf{Zheng \& Wei, \emph{J.~Chem.~Phys.}\
      \textbf{134}, 194101 (2011), doi:10.1063/1.3581031}. Both can be cited;
      only the second supports the PBNP-hybrid claim.
\end{enumerate}

\subsection*{Strongest novelty claims the writeup can defensibly make}

In order of strength:
\begin{enumerate}
\item \textbf{Application of PNP log-density ($u=\ln c$) into ORR--BV
      electrocatalysis.} No ORR/RDE/fuel-cell PNP--BV solver in the surveyed
      peer-reviewed literature uses $\ln c$ as primary Galerkin unknown.
\item \textbf{The log-rate BV construction itself, in a published PNP--BV
      solver.} The closest \emph{published} analogs are in different domains
      (viscoelastic fluids --- Fattal--Kupferman 2004; combustion --- CHEMKIN
      PLOG; ML --- log-sum-exp).
\item \textbf{PBNP applied specifically to ORR with Boltzmann-$\ClOf$} ---
      modest but unobjected.
\end{enumerate}
The first two are the load-bearing novelty claims; the third is a comfortable
``reuse in new setting.''

%% ============================================================
\section{Per-Change Literature Provenance}
%% ============================================================

\subsection{Change~\#1 --- PBNP Hybrid (3-species $+$ analytic Boltzmann counterion)}

\paragraph{Canonical PBNP paper.}
Zheng \& Wei (2011), \emph{J.~Chem.~Phys.}\ \textbf{134}(19), 194101,
\href{https://doi.org/10.1063/1.3581031}{doi:10.1063/1.3581031}, PMID~21599038,
PMC PMC3122111. Verified abstract: ``We propose an alternative model to reduce
number of Nernst--Planck equations to be solved \dots\ by substituting
Nernst--Planck equations with Boltzmann distributions of ion concentrations.''
This is exactly the writeup's split (3 active dynamic $+$ 1 analytic
Boltzmann).

\paragraph{Companion paper often confused with the above.}
Zheng, Chen, Wei (2011), \emph{J.~Comput.~Phys.}\ \textbf{230}(13),
5239--5262, \href{https://doi.org/10.1016/j.jcp.2011.03.020}{doi:10.1016/j.jcp.2011.03.020},
PMC PMC3087981. Full-PNP numerical-methods paper using
matched-interface-and-boundary methods --- \emph{not} PBNP. Cite as a high-order
full-PNP numerical reference; do not cite as a PBNP source.

\paragraph{Deeper conceptual ancestry.}
\begin{itemize}
\item Gouy 1910, \emph{J.~Phys.~Theor.~Appl.}\ \textbf{9}, 457--468 ---
      diffuse-layer model with Boltzmann ion statistics.
\item Chapman 1913, \emph{Phil.~Mag.}\ \textbf{25}, 475--481 --- closed-form
      Gouy--Chapman result.
\item Stern 1924, \emph{Z.~Elektrochem.}\ \textbf{30}, 508--516 --- combines
      Helmholtz $+$ Gouy--Chapman.
\end{itemize}
PBNP reduces to GCS as the active set vanishes and to full PNP as it expands.

\paragraph{Active-species-via-NP lineage (ion-channel biophysics).}
\begin{itemize}
\item Chen, Barcilon, Eisenberg 1992, \emph{Biophys.~J.}\ \textbf{61},
      1372--1393.
\item Chen, Eisenberg 1993, \emph{Biophys.~J.}\ \textbf{64}, 1405--1421,
      PMID~7686784.
\item Kurnikova, Coalson, Graf, Nitzan 1999, \emph{Biophys.~J.}\ \textbf{76},
      642--656.
\item Cardenas, Coalson, Kurnikova 2000, \emph{Biophys.~J.}\ \textbf{79},
      80--93, PMC1300917.
\item Lu, Holst, McCammon, Zhou 2010, \emph{J.~Comput.~Phys.}\ \textbf{229},
      6979--6994,
      \href{https://doi.org/10.1016/j.jcp.2010.05.035}{doi:10.1016/j.jcp.2010.05.035},
      PMC2922884 --- the most direct FEM-PNP predecessor that Wei's group
      built on.
\end{itemize}

\paragraph{Competitor reduction (worth knowing).}
Newman electroneutrality (Newman \& Thomas-Alyea, \emph{Electrochemical
Systems}, 3rd~ed., Wiley 2004; 4th~ed.\ with Balsara, 2019; Newman 1975,
\emph{AIChE J.}\ \textbf{21}, 25--41). Drops Poisson, enforces
$\sum_i z_i c_i=0$. Universal in fuel-cell engineering. Distinct from PBNP.

\paragraph{ORR-specific context.}
\begin{itemize}
\item Shinozaki et al.\ 2015, \emph{J.~Electrochem.~Soc.}\ \textbf{162},
      F1144--F1158,
      \href{https://doi.org/10.1149/2.1071509jes}{doi:10.1149/2.1071509jes}.
      Documents that $\ClOf$ is non-adsorbing/weakly adsorbing on Pt ---
      physical justification for treating it as Boltzmann-equilibrated.
\item Generalized modified PNP for PEMFC (\emph{Electrochim.~Acta} 2025,
      S0013468625004335): ongoing PNP-EDL modeling for fuel cells, but uses
      \emph{full} PNP, not PBNP.
\end{itemize}

\paragraph{Where the implementation aligns.}
Realizes Zheng--Wei's PBNP nearly exactly: 3 species
($\mathrm{O}_2,\HtwoOtwo,\Hpe$) with full Nernst--Planck transport plus a
Boltzmann analytic factor
$c_{\mathrm{ClO}_4}=c_{\mathrm{bulk}}\exp(-z_{\mathrm{ClO}_4}\varphi/\VT)$
with $z_{\mathrm{ClO}_4}=-1$, contributing to the Poisson residual. Confirmed
at \texttt{scripts/studies/v24\_3sp\_logc\_vs\_4sp\_validation.py:250-262}
and \texttt{scripts/studies/v18\_test\_3species\_boltzmann.py:109-147}.

%% ------------------------------------------------------------
\subsection{Change~\#2 --- Log-Concentration $u_i=\ln c_i$}

Three relevant lineages.

\paragraph{Lineage A --- Direct log-density / entropy-variable primary
(the writeup's path).}
\begin{itemize}
\item \textbf{Metti, Xu, Liu (2016)}, \emph{J.~Comput.~Phys.}\ \textbf{306},
      1--18,
      \href{https://doi.org/10.1016/j.jcp.2015.10.053}{doi:10.1016/j.jcp.2015.10.053}.
      Performs a logarithmic transformation of charge-carrier densities in FEM,
      proves a discrete energy estimate matching the continuous PNP energy
      law, enforces positivity. Canonical PNP log-density FEM paper. The phrase
      ``log-density formulation'' appears coined retroactively by Fu--Xu and
      successors; the technique is in the paper, the exact phrase may not be
      in the abstract.
\item \textbf{Fu \& Xu (2022)}, \emph{CMAME} \textbf{395}, 115031,
      \href{https://doi.org/10.1016/j.cma.2022.115031}{doi:10.1016/j.cma.2022.115031},
      arXiv:2105.01163. Discretizes the entropy variable
      $u_i=U'(c_i)=\log c_i$ directly with FEM in space $+$ DG in time;
      positivity by construction ($c_i=\exp(u_i)>0$); unconditional energy
      stability at arbitrary order. \textbf{This is the paper currently
      mis-cited as ``Liu, Maxian, Shen, Zitelli, Waltz, Monteiro 2022, CMAME
      396, 115070.''} That citation is fabricated.
\end{itemize}

\paragraph{Lineage B --- Slotboom variables / excess-chemical-potential flux
(semiconductor).}
\begin{itemize}
\item \textbf{Slotboom 1969}, \emph{Electronics Letters} \textbf{5}(26),
      677--678,
      \href{https://doi.org/10.1049/el:19690510}{doi:10.1049/el:19690510}.
      Canonical ``Slotboom variables'' reference: $\Phi_n=n\exp(-\psi/\VT)$.
\item \textbf{Slotboom 1973}, \emph{IEEE TED} \textbf{ED-20}, 669--679.
      Standard ``1973 Slotboom'' reference (bipolar transistor analysis).
\item \textbf{Slotboom 1977}, \emph{Solid-State Electron.}\ \textbf{20}(4),
      279--283. The actual ``PN-product in silicon'' paper ---
      band-gap-narrowing, \textbf{not} about log-density.
\item \textbf{Scharfetter \& Gummel 1969}, \emph{IEEE TED} \textbf{ED-16}(1),
      64--77. Exponential-fitting / upwind FV on $c$-primary; Bernoulli-function
      edge fluxes. \textbf{Not} a log-density substitution.
\item \textbf{Brezzi, Marini, Pietra 1989}, \emph{SIAM J.~Numer.~Anal.}\
      \textbf{26}(6), 1342--1355. Mixed-FE generalization of SG.
\item \textbf{Markowich 1986}, \emph{The Stationary Semiconductor Device
      Equations}, Springer,
      \href{https://doi.org/10.1007/978-3-7091-3678-2}{doi:10.1007/978-3-7091-3678-2}.
\item \textbf{Markowich, Ringhofer, Schmeiser 1990}, \emph{Semiconductor
      Equations}, Springer.
\item \textbf{Selberherr 1984}, \emph{Analysis and Simulation of Semiconductor
      Devices}, Springer.
\end{itemize}

\paragraph{Mapping between Slotboom and log-density.}
$\ln(\Phi_n)=\ln n - \psi/\VT$, so log-density and Slotboom differ only by an
affine map \emph{as continuous variables}. Galerkin discretizations on
$(\Phi_n,\psi)$ versus $(\ln n,\psi)$ yield different stiffness matrices.
The writeup must not equate them.

\paragraph{Lineage C --- Modern energy-stable / SAV / IEQ schemes for PNP.}

\begin{center}
\small
\begin{tabular}{@{}p{0.42\textwidth} p{0.52\textwidth}@{}}
\toprule
\textbf{Citation} & \textbf{What it does} \\
\midrule
H.~Liu \& Maimaitiyiming 2021, \emph{J.~Sci.~Comput.}\ \textbf{87}(3), 92 &
FD $c$-primary positivity $+$ energy stability for multi-D PNP \\
C.~Liu, Wang, Wise, Yue, Zhou, arXiv:2009.08076 (later \emph{Math.~Comp.}\ 2022) &
$H^{-1}$ gradient flow with singular log potential treated implicitly \\
Huang \& Shen 2021, \emph{SIAM J.~Sci.~Comput.}\ \textbf{43}(3), B746--B759,
doi:10.1137/20M1365417 &
SAV scheme; $c$-primary; bound-preservation via auxiliary variable \\
Canc\`es, Chainais-Hillairet, Gaudeul, Fuhrmann 2022,
\emph{Numer.~Math.}\ \textbf{151}, 949--986, doi:10.1007/s00211-022-01279-y &
Entropy-FV with size exclusion; modern bridge to electrochemistry \\
Farrell, Rotundo, Doan, Kantner, Fuhrmann, Koprucki book chapter (2017) &
Surveys flux discretizations including Slotboom and excess-chemical-potential \\
\bottomrule
\end{tabular}
\end{center}

\paragraph{Conceptual map.}
PNP free energy
$E[c,\psi]=\sum_i\!\int c_i(\ln c_i-1)\,\mathrm{d}x+\tfrac12\varepsilon\!\int|\nabla\psi|^2\,\mathrm{d}x-\sum_i z_i\!\int c_i\psi\,\mathrm{d}x$.
Three discretization strategies:
\begin{itemize}
\item[(A)] \textbf{Direct log-density / entropy-variable primary}
      (Metti--Xu--Liu $\to$ Fu--Xu).
\item[(B)] \textbf{Excess-chemical-potential flux} (Slotboom $\to$
      Brezzi--Marini--Pietra $\to$ Canc\`es).
\item[(C)] \textbf{SAV / IEQ} (Shen--Xu--Yang $\to$ Huang--Shen).
\end{itemize}
The writeup uses~(A).

\paragraph{Why electrochemistry diverged from semiconductor practice.}
Per Selberherr 1984 ch.~5 and Markowich--Ringhofer--Schmeiser 1990 ch.~3,
Slotboom variables work poorly under degenerate Fermi--Dirac statistics;
modern semiconductor codes prefer quasi-Fermi or excess-chemical-potential.
Electrochemistry rarely needs degeneracy correction so could use Slotboom,
but the field never adopted it --- codes descended from Newman's
CONDUC/DUALFOIL $c$-primary FD lineage rather than semiconductor TCAD.

\paragraph{Where the implementation aligns.}
\texttt{forms\_logc.py:55-56} constructs
\texttt{MixedFunctionSpace([V\_scalar]*(n+1))} with $u_i,\varphi$ all in
CG-order on the same space; \texttt{forms\_logc.py:177-178} confirms
$u_i=\ln c_i$ is the primary unknown and $c_i$ is a derived expression
$c_i=\exp(u_i)$ (with symmetric \texttt{\_U\_CLAMP=$\pm$30} overflow
protection --- \texttt{forms\_logc.py:185-198}). Flux uses the log-transform
identity $\nabla c=c\nabla u$ directly (\texttt{forms\_logc.py:281}). Time
stepping is backward Euler (\texttt{forms\_logc.py:283-285}), single past
state (\emph{not} BDF2).

%% ------------------------------------------------------------
\subsection{Change~\#3 --- Log-Rate Butler--Volmer}

The construction: rewrite
$r=k_0\,c_i\exp(-\alpha n_e\eta/\VT)$
as
$r=\exp(\ln k_0+u_i-\alpha n_e\eta/\VT)$ with $u_i=\ln c_i$, evaluated as one
exponentiation at the end. The cathodic branch folds in stoichiometric power
\emph{inside} the log:
\[
\ln r_{\mathrm{cat}}
\;=\; \ln k_0 + u_{\mathrm{cat}}
       + \sum_f \nu_f\bigl(u_{\mathrm{sp}(f)} - \ln c_{\mathrm{ref},f}\bigr)
       - \alpha\, n_e\,\eta/\VT.
\]

\paragraph{Tafel 1905 --- historical anchor, conceptually different.}
J.~Tafel, \emph{Z.~Phys.~Chem.}\ \textbf{50U}(1), 641--712,
\href{https://doi.org/10.1515/zpch-1905-5043}{doi:10.1515/zpch-1905-5043}.
Verified via De Gruyter, ESTIR Historic Papers in Electrochemistry, Russian
Mendeleev archive.

\paragraph{Critical conceptual distinction (must not conflate).}
\begin{itemize}
\item \textbf{Tafel 1905} $=$ high-overpotential physical \emph{approximation}
      that \textbf{drops one branch}: $\log|i|\approx \mathrm{const} + (\alpha
      nF/RT)\eta$.
\item \textbf{Log-rate BV} $=$ exact algebraic identity that \textbf{drops
      nothing}: $r=\exp(\ln k_0 + \ln c_i - \alpha n_e\eta/\VT)$.
\end{itemize}

\paragraph{Frumkin 1933 --- physical correction, also distinct.}
A.~N.~Frumkin, \emph{Z.~Phys.~Chem.~A} \textbf{164}, 121--133. Modifies the BV
potential drop: $\eta_{\mathrm{eff}}=\eta-\varphi_d$. Different correction
(interfacial physics), not numerical.

\paragraph{Bazant generalized BV and the gFBV$+$PNP framework --- physical
generalizations.}
\begin{itemize}
\item Bazant 2013, \emph{Acc.~Chem.~Res.}\ \textbf{46}, 1144--1160,
      \href{https://doi.org/10.1021/ar300145c}{doi:10.1021/ar300145c},
      PMID~23520980, arXiv:1208.1587. Generalizes BV via nonequilibrium
      thermodynamics; physical not numerical.
\item Bazant, Thornton, Ajdari 2004, \emph{Phys.~Rev.~E} \textbf{70}, 021506,
      \href{https://doi.org/10.1103/PhysRevE.70.021506}{doi:10.1103/PhysRevE.70.021506},
      PMID~15447495. PNP coupled to gFBV. \textbf{Standard exponential BV form
      throughout, no log-form rate evaluation.}
\item Bazant, Kilic, Storey, Ajdari 2009, \emph{Adv.~Colloid Interface Sci.}\
      \textbf{152}, 48--88,
      \href{https://doi.org/10.1016/j.cis.2009.10.001}{doi:10.1016/j.cis.2009.10.001}.
\item van Soestbergen, Biesheuvel, Bazant 2010, \emph{Phys.~Rev.~E}
      \textbf{81}, 021503.
\item van Soestbergen 2012, \emph{Russ.~J.~Electrochem.}\ \textbf{48},
      570--579,
      \href{https://doi.org/10.1134/S1023193512060110}{doi:10.1134/S1023193512060110}.
\end{itemize}
These represent the state of the art in PNP--BV. None address the
\emph{numerical evaluation} of BV as a stiff exponential.

\paragraph{Algebraic / inversion-direction reformulations (related but distinct
problem).}
\begin{itemize}
\item Khalili et al.\ 2016, ASME \emph{J.~Electrochem.~En.~Conv.~Stor.}\
      \textbf{13}(2), 021003. Algebraic reformulation to \emph{invert} BV for
      $\eta$.
\item Nasser \& Mantegazza 2022, \emph{J.~Phys.~Chem.~C},
      doi:10.1021/acs.jpcc.1c09620. Deformed exponentials for fitting.
\item COMSOL battery/electrochemistry: re-references overpotential against
      fixed activity reference; re-parametrization, not log-form.
\end{itemize}

\paragraph{The closest published structural cousin --- Fattal--Kupferman 2004.}
\begin{itemize}
\item \textbf{Fattal \& Kupferman 2004}, \emph{J.~Non-Newtonian Fluid Mech.}\
      \textbf{123}(2--3), 281--285,
      \href{https://doi.org/10.1016/j.jnnfm.2004.08.008}{doi:10.1016/j.jnnfm.2004.08.008}.
      Verified via Hebrew~U.\ CRIS and ScienceDirect. Citation count
      $\approx 484$. Key idea: ``transform a large class of differential
      constitutive models into an equation for the matrix logarithm of the
      conformation tensor''; high-Weissenberg numerical instability is
      ``the failure of polynomial-based approximations to properly represent
      exponential profiles.''
\item Companion: Fattal \& Kupferman 2005, \emph{J.~Non-Newtonian Fluid
      Mech.}\ \textbf{126}, 23--37,
      \href{https://doi.org/10.1016/j.jnnfm.2004.12.003}{doi:10.1016/j.jnnfm.2004.12.003}.
\end{itemize}

\paragraph{Exact mapping.}
\begin{itemize}
\item Fattal--Kupferman: $\Psi=\log(\mathrm{Conformation})$; discretize
      $\Psi$; $\exp(\Psi)$ only when needed.
\item Writeup: $\ln r=\ln k_0+u_i-\alpha n_e\eta/\VT$; assemble $\ln r$;
      $r=\exp(\ln r)$ once.
\end{itemize}

\paragraph{Stiff combustion / chemistry.}
CHEMKIN-II/III (Kee, Rupley, Miller, SAND-89-8009; SAND96-8216) --- PLOG
(pressure-dependent logarithmic interpolation) reactions. Lu \& Law 2009,
\emph{Prog.~Energy Combust.~Sci.}\ \textbf{35}, 192--215. Pope 1997,
\emph{Combust.~Theory Model.}\ \textbf{1}, 41--63. Higham et al.\ 2021,
\emph{IMA J.~Numer.~Anal.}\ \textbf{41}, 2311--2330,
\href{https://doi.org/10.1093/imanum/draa038}{doi:10.1093/imanum/draa038}
(log-sum-exp / softmax formal analysis).

The ``compute log of the rate, then exponentiate once'' pattern is well
known in CHEMKIN-style stiff chemistry codes and the closely related
log-sum-exp idiom in machine learning. \textbf{What is missing is the
explicit transplantation into electrochemistry / Butler--Volmer evaluation
as a published numerical strategy in a PNP--BV solver paper.}

\paragraph{Where the implementation aligns.}
Exact algebra at \texttt{Forward/bv\_solver/forms\_logc.py:324-356} (cathodic
324--339; anodic 340--356). The writeup's existing pointer ``see
\texttt{forms\_logc.py:294-360}'' undershoots; the full \texttt{if/else} block
runs lines 299--389. \textbf{Recommend pointer
\texttt{forms\_logc.py:294-388}.} Clamp-bypass argument documented at
\texttt{forms\_logc.py:294-298} and
\texttt{docs/CHATGPT\_HANDOFF\_6\_LOGRATE\_BREAKTHROUGH.md:107-118}.

%% ============================================================
\section{Per-Change Deviation Analysis}
%% ============================================================

\subsection{Change~\#1 --- Deviations from cited PBNP literature}

\paragraph{Reused (non-novel).}
The Boltzmann-passive $+$ NP-active reduction; sign convention
$z_{\mathrm{ClO}_4}=-1$; clamping $\varphi$ at $\pm 50$ dimensionless units
($\approx \pm\SI{1.28}{\volt}$) before exponentiation; non-dimensional
\texttt{charge\_rhs} prefactor.

\paragraph{In implementation but not in cited literature.}
\begin{enumerate}
\item \textbf{Application to ORR / \SI{0.1}{\mole\per\liter} $\mathrm{HClO}_4$.}
      No surveyed peer-reviewed paper Boltzmann-distributes $\ClOf$ (or any
      non-adsorbing anion) inside a PNP--BV ORR forward solver. The PEMFC
      PNP literature (Bessler, Eikerling, Weber--Newman) uses either full PNP
      or Newman electroneutrality. The construction is Zheng--Wei 2011; the
      application is the contribution.
\item \textbf{Architectural fragility --- \texttt{add\_boltzmann()}
      monkey-patch.} The Boltzmann contribution is \emph{not} in
      \texttt{Forward/bv\_solver/forms\_logc.py}. It is duplicated as a
      closure in 7 study scripts (\texttt{v18\_logc\_lsq\_inverse.py:234},
      \texttt{v24\_3sp\_logc\_vs\_4sp\_validation.py:250},
      \texttt{v18\_logc\_diagnostics.py:117},
      \texttt{v18\_logc\_noise\_sensitivity.py:85},
      \texttt{v19\_lograte\_extended\_adjoint\_check.py:166},
      \texttt{plot\_iv\_curves\_3sp\_true.py:108},
      \texttt{v18\_test\_3species\_boltzmann.py:109}). The bare
      \texttt{forms\_logc.py} Poisson residual contains only the
      three-explicit-species sum at line~418. Tightening recommendation, not
      literature finding.
\end{enumerate}

\paragraph{Genuine novelty assessment.}
Modest. Construction fully reused; only the application domain (ORR PNP--BV
with Boltzmann $\ClOf$) appears undocumented.

\subsection{Change~\#2 --- Deviations from cited log-density literature}

\paragraph{Reused (non-novel).}
\begin{enumerate}
\item The substitution $u_i=\ln c_i$, recovery $c_i=\exp(u_i)$, flux identity
      $\nabla c=c\,\nabla u$ --- directly from Metti--Xu--Liu 2016 and Fu--Xu
      2022.
\item Energy-stable interpretation: $c_i=\exp(u_i)>0$ $\Rightarrow$ positivity
      by construction; discrete entropy / energy stability arguments transfer.
\item Monolithic Newton over Gummel splitting --- standard in modern
      energy-stable PNP literature.
\item Same CG order on $u_i$ and $\varphi$ --- standard in Metti--Xu--Liu and
      FEM-PNP electrochemistry generally.
\end{enumerate}

\paragraph{In implementation but not in cited literature.}
\begin{enumerate}
\item \textbf{PNP--BV ORR application} (load-bearing novelty). After targeted
      search, no peer-reviewed electrochemistry-side ORR / RDE / fuel-cell /
      battery PNP--BV solver paper uses $\ln c$ as a primary Galerkin unknown.
      Standard practice in that community is $c$-primary with adaptive
      meshing. The closest electrochemistry-side log-density work is in Li-ion
      concentrated-solution theory (Newman group, Doyle--Fuller--Newman
      1993$+$), where chemical potential $\mu=\mu_0+RT\ln(\gamma c)$ appears
      in fluxes --- but they still discretize $c$, not $\ln c$.
\item \textbf{Propagation of $u_i$ into the BV boundary condition}
      (Change~\#3). Metti--Xu--Liu and Fu--Xu handle bulk PNP transport only.
\item \textbf{Symmetric \texttt{\_U\_CLAMP=$\pm$30} as overflow protection
      rather than concentration floor} (\texttt{forms\_logc.py:185-198}).
      Comment: ``$\exp(\pm 30)$ covers $[9.4\!\times\!10^{-14},\,1.07\!\times\!10^{13}]$,
      adequate for typical EDL profiles.'' Bypassed in the log-rate boundary
      path. Small construction choice, not in any specific cited paper,
      consistent with general floating-point hygiene.
\item \textbf{\texttt{\_C\_FLOOR=1e-20} Dirichlet-BC value for product
      species} (\texttt{forms\_logc.py:431-437}) to avoid $\ln 0$ for $\HtwoOtwo$
      initial concentration. Small local construction.
\end{enumerate}

\paragraph{Genuine novelty assessment.}
Genuine --- but as a \emph{domain transfer}, not a new substitution. The
substitution itself is old (semiconductor 1969, mathematical PNP 2016). No
surveyed ORR--BV PNP solver uses it. Confidence medium-high (proving a
negative is impossible, but targeted search returned nothing).

\subsection{Change~\#3 --- Deviations from cited / surveyed BV literature}

\paragraph{Reused (non-novel).}
\begin{enumerate}
\item Standard BV cathodic kinetic expression
      $r=k_0 c_i \exp(-\alpha n_e\eta/\VT)$ --- textbook Bockris--Reddy /
      Bard--Faulkner.
\item Stoichiometric-power loop in cathodic branch ---
      $\prod_f c_{\mathrm{sp}(f)}^{\nu_f}$ --- textbook mass-action.
\item The algebraic identity itself ---
      $r=k_0 c_i \exp(\dots)\Leftrightarrow\exp(\ln k_0+\ln c_i+\dots)$ --- is
      trivial. Any electrochemist could derive it in two lines.
\item ``Evaluate inside the log to avoid overflow at extreme rates'' pattern
      --- well established in stiff combustion (CHEMKIN PLOG), viscoelastic
      flow (Fattal--Kupferman), and ML (log-sum-exp). Standard machinery.
\end{enumerate}

\paragraph{In implementation but not in surveyed literature.}
\begin{enumerate}
\item \textbf{The literal log-rate BV identity used as a numerical strategy in
      a PNP--BV / ORR solver.} No electrochemistry paper located that
      publishes this construction in a PNP--BV solver context. The closest
      published cousin is Fattal--Kupferman 2004 in viscoelastic fluid
      mechanics --- structurally identical, domain different. The ``small
      local numerical novelty'' hedge is defensible.
\item \textbf{Clamp-bypass as the primary motivation.} The implementation's
      reason is specifically: ``evaluating BV against \texttt{\_U\_CLAMP}-bounded
      $c_{\mathrm{surf}}$ produces a phantom $\mathrm{R}_2$ cathodic sink at
      low surface $\HtwoOtwo$.'' Mechanism per
      \texttt{docs/CHATGPT\_HANDOFF\_6\_LOGRATE\_BREAKTHROUGH.md:107-118}:
      when Newton needs $c_{\HtwoOtwo}$ below $\exp(-30)\approx 9.4\!\times\!10^{-14}$,
      the symmetric clamp pins it; combined with a saturated $\exp(50)$ in
      $\mathrm{R}_2$'s BV exponent, the floor times the huge exponential
      gives a spurious $\mathrm{R}_2$ sink that nothing else can balance.
      Log-rate evaluates
      $\exp(\ln k_0+u_{\HtwoOtwo}+2(u_\mathrm{H}-\ln c_{\mathrm{ref}})-\alpha n_e\eta)$
      so $u_{\HtwoOtwo}$ enters additively and can be arbitrarily negative;
      the expression decays smoothly to zero. The specific clamp-bypass
      argument is not in any cited literature --- it is a property of
      \emph{this} implementation's interaction.
\item \textbf{Asymmetry: cathodic branch carries the full sum loop; anodic
      branch is single-species} (\texttt{forms\_logc.py:340-356}). The
      writeup's ``both branches subtracted with
      $\sum_f \nu_f(u_{\mathrm{sp}(f)}-\ln c_{\mathrm{ref},f})$'' is
      \emph{partially} accurate: cathodic has the full loop (lines 331--337);
      anodic has only the bare anodic-species term $u_{\mathrm{anod}}$
      (lines 343--345 main case, 350--353 fallback). For production,
      $\mathrm{R}_2$ is irreversible
      (\texttt{v24\_3sp\_logc\_vs\_4sp\_validation.py:233}), so the anodic
      branch falls into else at line 355 $\to$ $\mathrm{anodic}=0$.
      Algebra-description should be tightened.
\item \textbf{Line-range pointer drift.} Writeup says
      \texttt{forms\_logc.py:294-360}. The codebase shows the log-rate true
      branch at 324--356; full \texttt{if/else} block runs 294--388.
      \textbf{Recommend \texttt{forms\_logc.py:294-388}.}
\end{enumerate}

\paragraph{Genuine novelty assessment.}
Small local numerical novelty (defensible). Construction mathematically
trivial; deployment in a PNP--BV / ORR solver undocumented in the surveyed
literature; structurally identical machinery exists in different fields. The
unique angle is the clamp-bypass mechanism specific to this code's
\texttt{\_U\_CLAMP}-vs-saturated-BV interaction.

%% ============================================================
\section{Surrounding-Methods Provenance}
%% ============================================================

\subsection{ORR Mechanism / RRDE}

\paragraph{Two-step equilibrium potentials.}
$\mathrm{R}_1$: $\mathrm{O}_2+2\Hpe+2e^-\to\HtwoOtwo$,
$\Eeq^{(1)}=+\SI{0.695}{\volt}$ vs SHE ($\approx \SI{0.68}{\volt}$).
$\mathrm{R}_2$: $\HtwoOtwo+2\Hpe+2e^-\to 2\mathrm{H}_2\mathrm{O}$,
$\Eeq^{(2)}=+\SI{1.776}{\volt}$ vs SHE ($\approx \SI{1.78}{\volt}$). Sum: 4-$e^-$
ORR with $\Eeq=+\SI{1.229}{\volt}$ vs SHE. Cite Bard, Faulkner \& White 3rd~ed.\
(2022) ch.~9.

Production reaction config
(\texttt{v24\_3sp\_logc\_vs\_4sp\_validation.py:222-237}): $\mathrm{R}_1$ has
\texttt{E\_eq\_v=0.68}, \texttt{n\_electrons=2}, cathodic $\mathrm{O}_2$,
anodic $\HtwoOtwo$, stoichiometry $[-1,+1,-2]$ with $\Hpe$ power 2;
$\mathrm{R}_2$ has \texttt{E\_eq\_v=1.78}, \texttt{n\_electrons=2},
irreversible.

\paragraph{Foundational mechanism papers.}
\begin{itemize}
\item Wroblowa, Yen-Chi-Pan, Razumney 1976,
      \emph{J.~Electroanal.~Chem.}\ \textbf{69}(2), 195--201,
      doi:10.1016/S0022-0728(76)80250-1. THE classic 2-step ORR paper.
\item Damjanovic, Genshaw, Bockris 1967,
      \emph{J.~Electrochem.~Soc.}\ \textbf{114}(5), 466--472. RRDE
      $\HtwoOtwo$-on-Pt mechanism.
\item Markovi\'c, Schmidt, Stamenkovi\'c, Ross 2001, \emph{Fuel Cells}
      \textbf{1}(2), 105--116.
\item Markovi\'c \& Ross 2002, \emph{Surf.~Sci.~Rep.}\ \textbf{45}(4--6),
      117--229.
\item N{\o}rskov, Rossmeisl et al.\ 2004, \emph{J.~Phys.~Chem.~B}
      \textbf{108}(46), 17886--17892, doi:10.1021/jp047349j.
\item Stamenkovi\'c et al.\ 2007, \emph{Science} \textbf{315}(5811), 493--497,
      doi:10.1126/science.1135941.
\item Kulkarni, Siahrostami, Patel, N{\o}rskov 2018, \emph{Chem.~Rev.}\
      \textbf{118}(5), 2302--2312, doi:10.1021/acs.chemrev.7b00488.
\end{itemize}

\subsection{Kinetic Inverse Problems and Tafel-Ridge Identifiability}

\paragraph{Battery / Newman-style PDE-constrained ID.}
\begin{itemize}
\item Bizeray, Kim, Duncan, Howey 2017/2018,
      \emph{IEEE Trans.~Control Syst.~Technol.}\ (preprint
      arXiv:1702.02471). SPM Li-ion has only six identifiable parameter
      groups.
\item Forman, Moura, Stein, Fathy 2012, \emph{J.~Power Sources} \textbf{210},
      263--275, doi:10.1016/j.jpowsour.2012.03.009. DFN parameter ID $+$
      post-hoc FIM.
\item Laue, R\"oder, Krewer 2021, \emph{J.~Appl.~Electrochem.}\ \textbf{51},
      1253--1265, doi:10.1007/s10800-021-01579-5.
\item Park, Kato, Gima, Klein, Moura 2018, \emph{J.~Electrochem.~Soc.}\
      \textbf{165}(7), A1309. Convex OED for SPM via FIM.
\end{itemize}

\paragraph{Identifiability theory.}
\begin{itemize}
\item Raue et al.\ 2009, \emph{Bioinformatics} \textbf{25}(15), 1923--1929,
      doi:10.1093/bioinformatics/btp358. Profile likelihood; structural vs
      practical non-identifiability cleavage.
\item Brun, Reichert, K\"unsch 2001, \emph{Water Resour.~Res.}\
      \textbf{37}(4), 1015--1030, doi:10.1029/2000WR900350. Sensitivity index
      $+$ collinearity index $=1/\sigma_{\min}(J_{\mathrm{normalized}})$.
\item Quaiser \& M\"onnigmann 2009, \emph{BMC Syst.~Biol.}\ \textbf{3}, 50,
      doi:10.1186/1752-0509-3-50. Eigenvalue method.
\end{itemize}

\paragraph{Tafel-ridge directly.}
In the BV Tafel branch, $\log|i|\approx \log i_0 + (1-\alpha) f \eta$:
$\log i_0$ enters additively, $\alpha$ as slope; both identifiable. The ridge
emerges when (a)~the $\eta$ window is narrow or (b)~mass transport /
$\mathrm{R}_2$ coupling bends the apparent Tafel slope.
\begin{itemize}
\item Schalenbach et al.\ 2024, \emph{ACS Energy Lett.},
      doi:10.1021/acsenergylett.4c00266. Variance in $(\alpha,i_0)$ across
      groups.
\item Anantharaj \& Noda 2021, \emph{Sci.~Rep.}\ \textbf{11}, 8915,
      doi:10.1038/s41598-021-87951-z. Differential Tafel decoupler.
\item Bilodeau, Gibson, Garnett 2021, \emph{Nat.~Commun.}\ \textbf{12}, 825,
      doi:10.1038/s41467-021-20924-y. Bayesian $(i_0,\alpha)$ elongated
      posterior in $\mathrm{CO}_2\mathrm{RR}$.
\end{itemize}

\subsection{Charge / Voltage Continuation}

\begin{itemize}
\item Keller 1977, ``Numerical solution of bifurcation and nonlinear
      eigenvalue problems,'' in \emph{Applications of Bifurcation Theory}
      (Rabinowitz, ed.), Academic Press. Pseudo-arclength continuation.
\item Allgower \& Georg 1990/2003 SIAM Classics, \emph{Numerical Continuation
      Methods}.
\item Doedel et al.\ AUTO-07P / AUTO-2000.
\item Uecker 2021, \emph{Jahresber.~DMV} \textbf{123}, 199--248,
      doi:10.1365/s13291-021-00241-5.
\item Markowich, Ringhofer, Schmeiser 1990, \emph{Semiconductor Equations},
      Springer, doi:10.1007/978-3-7091-6961-2. \textbf{Textbook origin of
      bias-ramping for drift-diffusion / PNP.}
\item Hyon, Eisenberg, Liu 2010, \emph{Commun.~Math.~Sci.}\ \textbf{9}(2),
      459--475. Energetic-variational PNP with steric.
\item Kilic, Bazant, Ajdari 2007, \emph{Phys.~Rev.~E} \textbf{75}, 021502
      (I), 021503 (II), doi:10.1103/PhysRevE.75.021502. Why classical PB/PNP
      fails $>\SI{25}{\milli\volt}$.
\item Schmuck, Bazant 2015, \emph{SIAM J.~Appl.~Math.}\ \textbf{75}(3),
      1369--1401, doi:10.1137/140968082.
\item EchemFEM (Roy, Lin, Hahn 2024), JOSS / LLNL-JRNL-860653. Most direct
      contemporary Firedrake-PNP reference.
\end{itemize}

\subsection{Adjoint / Discrete-Adjoint in Firedrake}

\begin{itemize}
\item Farrell, Ham, Funke, Rognes 2013, \emph{SIAM J.~Sci.~Comput.}\
      \textbf{35}(4), C369--C393, doi:10.1137/120873558. Algorithmic basis of
      \texttt{dolfin-adjoint}.
\item Mitusch, Funke, Dokken 2019, \emph{J.~Open Source Softw.}\ \textbf{4}(38),
      1292, doi:10.21105/joss.01292. \texttt{pyadjoint} software paper.
\item Plessix 2006, \emph{Geophys.~J.~Int.}\ \textbf{167}(2), 495--503,
      doi:10.1111/j.1365-246X.2006.02978.x. Adjoint-state method review.
\item Gunzburger 2003, \emph{Perspectives in Flow Control and Optimization},
      SIAM.
\item Hinze, Pinnau, Ulbrich, Ulbrich 2009, \emph{Optimization with PDE
      Constraints}, Springer.
\item Ascher \& Petzold 1998, \emph{Computer Methods for ODEs and DAEs}, SIAM.
\end{itemize}
Taylor-test verification documented in Farrell--Ham--Funke--Rognes 2013, \S5.

\subsection{Local Fisher Information \& Canonical-Ridge Cosines}

\begin{itemize}
\item Fisher 1922, \emph{Phil.~Trans.~R.~Soc.~A} \textbf{222}, 309--368,
      doi:10.1098/rsta.1922.0009.
\item Lehmann \& Casella 1998, \emph{Theory of Point Estimation}, 2nd~ed.,
      Springer.
\item Brun, Reichert, K\"unsch 2001 (above). Their collinearity index is
      $1/\sigma_{\min}(J_{\mathrm{normalized}})$.
\item Quaiser, M\"onnigmann 2009 (above). Eigenvalue method $=$ SVD of the FIM.
\item Pant 2018, \emph{J.~R.~Soc.~Interface} \textbf{15}, 20170871,
      doi:10.1098/rsif.2017.0871. Most modern $v_{\min}(\mathrm{FIM})$
      inspection.
\item Hotelling 1936 (\emph{Biometrika}), canonical correlation analysis.
\item van der Vaart 2000, \emph{Asymptotic Statistics}, \S5.
\item Vajda \& Rabitz 1989, \emph{J.~Phys.~Chem.}\ \textbf{93}, 5043. Overlap
      angles between sensitivity-derived eigenvectors.
\item Komorowski, Costa, Rand, Stumpf 2011, \emph{PNAS} \textbf{108}(21),
      8645--8650, doi:10.1073/pnas.1015814108.
\end{itemize}
The canonical-ridge cosine is the user's naming for the dot-product of
$v_{\min}(\mathrm{FIM})$ with a chosen analytical Tafel-ridge direction. Pant
2018 is the closest pre-existing diagnostic.

\textbf{``Joshi, Seidel-Morgenstern, Tiller 2006''} --- not located in any
database. Drop or replace with Brun 2001 $+$ Quaiser 2009 $+$ Pant 2018.

\subsection{Monolithic vs Splitting PNP Solvers}

\begin{itemize}
\item Gummel 1964, \emph{IEEE TED} \textbf{11}(10), 455--465. Original Gummel
      iteration.
\item Markowich, Ringhofer, Schmeiser 1990 (above), ch.~7.
\item Yu, Holst, McCammon 2007, \emph{J.~Comput.~Chem.}\ \textbf{28},
      1827--1839. Monolithic PNP, biological geometries.
\item Liu \& Wang 2014, \emph{J.~Comput.~Phys.}\ \textbf{268}, 363--376.
      Energy-stable structure-preserving FD.
\item EchemFEM (Roy, Lin, Hahn 2024).
\end{itemize}
The writeup uses monolithic Newton on the mixed
$(u_1,\dots,u_n,\varphi)$ space.

%% ============================================================
\section{Annotated Bibliography (verified)}
%% ============================================================

\subsection{Log-density / entropy-variable PNP}

\begingroup
\small
\setlength{\tabcolsep}{4pt}
\renewcommand{\arraystretch}{1.2}
\begin{longtable}{@{}p{0.46\textwidth} p{0.13\textwidth} p{0.36\textwidth}@{}}
\toprule
\textbf{Citation} & \textbf{Status} & \textbf{One-sentence summary} \\
\midrule
\endhead
Slotboom 1969, \emph{Electronics Letters} \textbf{5}(26), 677--678,
doi:10.1049/el:19690510 &
Verified &
Canonical Slotboom-variables paper; introduces $\Phi_n=n\exp(-\psi/\VT)$ \\
Slotboom 1973, \emph{IEEE TED} \textbf{ED-20}, 669--679 &
Verified &
Bipolar transistor analysis; not the PN-product paper \\
Slotboom 1977, \emph{Solid-State Electron.}\ \textbf{20}(4), 279--283 &
Verified &
Band-gap-narrowing paper, not log-density \\
Scharfetter \& Gummel 1969, \emph{IEEE TED} \textbf{ED-16}(1), 64--77 &
Verified &
Exponential-fitting / upwind FV on $c$-primary; \emph{not} log-density
substitution \\
Brezzi, Marini, Pietra 1989, \emph{SIAM J.~Numer.~Anal.}\ \textbf{26}(6),
1342--1355 &
Verified &
Mixed-FE generalization of SG \\
Markowich 1986, \emph{The Stationary Semiconductor Device Equations}, Springer,
doi:10.1007/978-3-7091-3678-2 &
Verified &
Surveys variable choices \\
Markowich, Ringhofer, Schmeiser 1990, \emph{Semiconductor Equations}, Springer
&
Verified &
Drift-diffusion $+$ iterative schemes textbook \\
Selberherr 1984, \emph{Analysis and Simulation of Semiconductor Devices},
Springer &
Verified &
Industrial Slotboom $+$ SG $+$ Gummel \\
Metti, Xu, Liu 2016, \emph{J.~Comput.~Phys.}\ \textbf{306}, 1--18,
doi:10.1016/j.jcp.2015.10.053 &
Verified &
Canonical PNP log-density FEM paper \\
Fu \& Xu 2022, \emph{CMAME} \textbf{395}, 115031,
doi:10.1016/j.cma.2022.115031, arXiv:2105.01163 &
Verified &
High-order space-time FEM for PNP via the entropy variable; \textbf{the paper
currently mis-cited as ``Liu, Maxian, Shen, Zitelli, Waltz, Monteiro 2022''} \\
H.~Liu \& Maimaitiyiming 2021, \emph{J.~Sci.~Comput.}\ \textbf{87}(3), 92,
doi:10.1007/s10915-021-01503-1 &
Verified &
FD $c$-primary positivity $+$ energy stability for multi-D PNP \\
C.~Liu, Wang, Wise, Yue, Zhou, arXiv:2009.08076 &
Verified &
$H^{-1}$ gradient flow with singular log potential \\
Huang \& Shen 2021, \emph{SIAM J.~Sci.~Comput.}\ \textbf{43}(3), B746--B759,
doi:10.1137/20M1365417 &
Verified &
SAV scheme for PNP/Keller--Segel \\
Canc\`es, Chainais-Hillairet, Gaudeul, Fuhrmann 2022, \emph{Numer.~Math.}\
\textbf{151}, 949--986, doi:10.1007/s00211-022-01279-y &
Verified &
Entropy-FV for PNP with size exclusion \\
\bottomrule
\end{longtable}
\endgroup

\subsection{PBNP hybrid and ion-channel PNP}

\begingroup
\small
\setlength{\tabcolsep}{4pt}
\renewcommand{\arraystretch}{1.2}
\begin{longtable}{@{}p{0.50\textwidth} p{0.13\textwidth} p{0.32\textwidth}@{}}
\toprule
\textbf{Citation} & \textbf{Status} & \textbf{Summary} \\
\midrule
\endhead
Zheng \& Wei 2011, \emph{J.~Chem.~Phys.}\ \textbf{134}(19), 194101,
doi:10.1063/1.3581031, PMID 21599038, PMC PMC3122111 &
Verified &
Canonical PBNP paper \\
Zheng, Chen, Wei 2011, \emph{J.~Comput.~Phys.}\ \textbf{230}(13), 5239--5262,
doi:10.1016/j.jcp.2011.03.020, PMC PMC3087981 &
Verified &
Second-order full-PNP solver, \emph{not} PBNP \\
Gouy 1910, \emph{J.~Phys.~Theor.~Appl.}\ \textbf{9}, 457--468 &
Verified &
Diffuse-layer model, Boltzmann ion statistics \\
Chapman 1913, \emph{Phil.~Mag.}\ \textbf{25}, 475--481 &
Verified &
Closed-form Gouy--Chapman \\
Stern 1924, \emph{Z.~Elektrochem.}\ \textbf{30}, 508--516 &
Verified &
Combines Helmholtz $+$ GC \\
Chen, Barcilon, Eisenberg 1992, \emph{Biophys.~J.}\ \textbf{61}, 1372--1393 &
Verified &
Open ionic channel PNP \\
Chen, Eisenberg 1993, \emph{Biophys.~J.}\ \textbf{64}, 1405--1421,
PMID 7686784 &
Verified &
Foundational PNP for one-conformation channels \\
Kurnikova, Coalson, Graf, Nitzan 1999, \emph{Biophys.~J.}\ \textbf{76},
642--656 &
Verified &
3D PNP for gramicidin A \\
Cardenas, Coalson, Kurnikova 2000, \emph{Biophys.~J.}\ \textbf{79}, 80--93,
PMC1300917 &
Verified &
3D PNP gramicidin A \\
Lu, Holst, McCammon, Zhou 2010, \emph{J.~Comput.~Phys.}\ \textbf{229},
6979--6994, doi:10.1016/j.jcp.2010.05.035, PMC2922884 &
Verified &
FEM-PNP for biomolecular diffusion-reaction \\
Newman \& Thomas-Alyea 2004, \emph{Electrochemical Systems}, 3rd~ed., Wiley,
ISBN 0-471-47756-7 &
Verified &
Electroneutrality reference \\
Newman \& Balsara 2019, \emph{Electrochemical Systems}, 4th~ed., Wiley-ECS &
Verified &
Current edition \\
Newman 1975, \emph{AIChE~J.}\ \textbf{21}, 25--41 &
Verified &
Porous-electrode theory \\
Shinozaki et al.\ 2015, \emph{J.~Electrochem.~Soc.}\ \textbf{162}, F1144--F1158,
doi:10.1149/2.1071509jes &
Verified &
RDE/ORR $\mathrm{HClO}_4$ benchmarking \\
\bottomrule
\end{longtable}
\endgroup

\subsection{Butler--Volmer, Tafel, Frumkin, generalized BV}

\begingroup
\small
\setlength{\tabcolsep}{4pt}
\renewcommand{\arraystretch}{1.2}
\begin{longtable}{@{}p{0.50\textwidth} p{0.18\textwidth} p{0.27\textwidth}@{}}
\toprule
\textbf{Citation} & \textbf{Status} & \textbf{Summary} \\
\midrule
\endhead
Tafel 1905, \emph{Z.~Phys.~Chem.}\ \textbf{50U}(1), 641--712,
doi:10.1515/zpch-1905-5043 &
Verified &
Original empirical Tafel law \\
Frumkin 1933, \emph{Z.~Phys.~Chem.~A} \textbf{164}, 121--133 &
Verified &
Frumkin correction $\eta_{\mathrm{eff}}=\eta-\varphi_d$ \\
Bazant, Thornton, Ajdari 2004, \emph{Phys.~Rev.~E} \textbf{70}, 021506,
doi:10.1103/PhysRevE.70.021506, PMID 15447495 &
Verified &
Diffuse-charge dynamics; PNP $+$ gFBV \\
Kilic, Bazant, Ajdari 2007, \emph{Phys.~Rev.~E} \textbf{75}, 021502/021503,
doi:10.1103/PhysRevE.75.021502 &
Verified &
Steric effects; classical PB fails $>\SI{25}{\milli\volt}$ \\
Bazant, Kilic, Storey, Ajdari 2009, \emph{Adv.~Colloid Interface Sci.}\
\textbf{152}, 48--88, doi:10.1016/j.cis.2009.10.001 &
Verified &
Induced-charge electrokinetics \\
van Soestbergen, Biesheuvel, Bazant 2010, \emph{Phys.~Rev.~E} \textbf{81},
021503 &
Verified &
PNP $+$ gFBV transient response \\
van Soestbergen 2012, \emph{Russ.~J.~Electrochem.}\ \textbf{48}, 570--579,
doi:10.1134/S1023193512060110 &
Verified &
Frumkin-BV $+$ mass transfer \\
Bazant 2013, \emph{Acc.~Chem.~Res.}\ \textbf{46}, 1144--1160,
doi:10.1021/ar300145c, PMID 23520980, arXiv:1208.1587 &
Verified &
Generalized BV via nonequilibrium thermodynamics \\
Bockris, Reddy, Gamboa-Aldeco 2002, \emph{Modern Electrochemistry} \textbf{2A},
2nd~ed., Kluwer/Plenum &
Verified (IA) &
Canonical BV/Tafel textbook \\
Bard, Faulkner, White 2022, \emph{Electrochemical Methods}, 3rd~ed., Wiley &
Verified &
Canonical RRDE/Levich/BV textbook \\
Khalili et al.\ 2016, ASME \emph{J.~Electrochem.~En.~Conv.~Stor.}\
\textbf{13}(2), 021003 &
Verified &
Algebraic BV reformulation (inverts for $\eta$) \\
\bottomrule
\end{longtable}
\endgroup

\subsection{Closest log-rate-evaluation cousins (different domains)}

\begingroup
\small
\setlength{\tabcolsep}{4pt}
\renewcommand{\arraystretch}{1.2}
\begin{longtable}{@{}p{0.46\textwidth} p{0.20\textwidth} p{0.30\textwidth}@{}}
\toprule
\textbf{Citation} & \textbf{Status} & \textbf{Summary} \\
\midrule
\endhead
\textbf{Fattal \& Kupferman 2004}, \emph{J.~Non-Newtonian Fluid Mech.}\
\textbf{123}(2--3), 281--285, doi:10.1016/j.jnnfm.2004.08.008 &
Verified &
Log-conformation tensor at high Weissenberg; \textbf{closest published
structural cousin to log-rate BV} \\
Fattal \& Kupferman 2005, \emph{J.~Non-Newtonian Fluid Mech.}\ \textbf{126},
23--37, doi:10.1016/j.jnnfm.2004.12.003 &
Verified &
Application paper for log-conformation \\
Kee, Rupley, Miller (CHEMKIN-II SAND-89-8009; CHEMKIN-III SAND96-8216) &
Verified &
PLOG (logarithmic interpolation) reactions \\
Lu \& Law 2009, \emph{Prog.~Energy Combust.~Sci.}\ \textbf{35}, 192--215 &
Verified (metadata) &
Stiff-chemistry numerics review \\
Pope 1997, \emph{Combust.~Theory Model.}\ \textbf{1}, 41--63 &
Verified (metadata) &
Tabulates log of rates \\
Higham et al.\ 2021, \emph{IMA J.~Numer.~Anal.}\ \textbf{41}, 2311--2330,
doi:10.1093/imanum/draa038 &
Verified &
Log-sum-exp / softmax numerical analysis \\
\bottomrule
\end{longtable}
\endgroup

\subsection{ORR mechanism}

\begin{center}
\small
\begin{tabular}{@{}p{0.78\textwidth} p{0.16\textwidth}@{}}
\toprule
\textbf{Citation} & \textbf{Status} \\
\midrule
Wroblowa, Yen-Chi-Pan, Razumney 1976,
\emph{J.~Electroanal.~Chem.~Interfacial Electrochem.}\ \textbf{69}(2),
195--201, doi:10.1016/S0022-0728(76)80250-1 &
Verified \\
Damjanovic, Genshaw, Bockris 1967, \emph{J.~Electrochem.~Soc.}\ \textbf{114}(5),
466--472 &
Verified \\
Markovi\'c, Schmidt, Stamenkovi\'c, Ross 2001, \emph{Fuel Cells} \textbf{1}(2),
105--116 &
Verified \\
Markovi\'c \& Ross 2002, \emph{Surf.~Sci.~Rep.}\ \textbf{45}(4--6), 117--229 &
Verified \\
N{\o}rskov, Rossmeisl et al.\ 2004, \emph{J.~Phys.~Chem.~B} \textbf{108}(46),
17886--17892, doi:10.1021/jp047349j &
Verified \\
Stamenkovi\'c et al.\ 2007, \emph{Science} \textbf{315}(5811), 493--497,
doi:10.1126/science.1135941 &
Verified \\
Kulkarni, Siahrostami, Patel, N{\o}rskov 2018, \emph{Chem.~Rev.}\
\textbf{118}(5), 2302--2312, doi:10.1021/acs.chemrev.7b00488 &
Verified \\
\bottomrule
\end{tabular}
\end{center}

\subsection{Inverse problems / identifiability / FIM}

\begingroup
\small
\setlength{\tabcolsep}{4pt}
\renewcommand{\arraystretch}{1.2}
\begin{longtable}{@{}p{0.78\textwidth} p{0.16\textwidth}@{}}
\toprule
\textbf{Citation} & \textbf{Status} \\
\midrule
\endhead
Fisher 1922, \emph{Phil.~Trans.~R.~Soc.~A} \textbf{222}, 309--368,
doi:10.1098/rsta.1922.0009 & Verified \\
Lehmann \& Casella 1998, \emph{Theory of Point Estimation}, 2nd~ed., Springer
& textbook \\
Brun, Reichert, K\"unsch 2001, \emph{Water Resour.~Res.}\ \textbf{37}(4),
1015--1030, doi:10.1029/2000WR900350 & Verified \\
Quaiser \& M\"onnigmann 2009, \emph{BMC Syst.~Biol.}\ \textbf{3}, 50,
doi:10.1186/1752-0509-3-50 & Verified \\
Raue et al.\ 2009, \emph{Bioinformatics} \textbf{25}(15), 1923--1929,
doi:10.1093/bioinformatics/btp358 & Verified \\
Pant 2018, \emph{J.~R.~Soc.~Interface} \textbf{15}, 20170871,
doi:10.1098/rsif.2017.0871 & Verified \\
Komorowski, Costa, Rand, Stumpf 2011, \emph{PNAS} \textbf{108}(21), 8645--8650,
doi:10.1073/pnas.1015814108 & Verified \\
Vajda \& Rabitz 1989, \emph{J.~Phys.~Chem.}\ \textbf{93}, 5043 & Verified \\
Bizeray, Kim, Duncan, Howey 2018, \emph{IEEE Trans.~Control Syst.~Technol.}\
(arXiv:1702.02471) & Verified \\
Forman, Moura, Stein, Fathy 2012, \emph{J.~Power Sources} \textbf{210},
263--275, doi:10.1016/j.jpowsour.2012.03.009 & Verified \\
Laue, R\"oder, Krewer 2021, \emph{J.~Appl.~Electrochem.}\ \textbf{51},
1253--1265, doi:10.1007/s10800-021-01579-5 & Verified \\
Park, Kato, Gima, Klein, Moura 2018, \emph{J.~Electrochem.~Soc.}\
\textbf{165}(7), A1309 & Verified \\
Bilodeau, Gibson, Garnett 2021, \emph{Nat.~Commun.}\ \textbf{12}, 825,
doi:10.1038/s41467-021-20924-y & Verified \\
Anantharaj \& Noda 2021, \emph{Sci.~Rep.}\ \textbf{11}, 8915,
doi:10.1038/s41598-021-87951-z & Verified \\
Schalenbach et al.\ 2024, \emph{ACS Energy Lett.},
doi:10.1021/acsenergylett.4c00266 & Verified \\
\bottomrule
\end{longtable}
\endgroup

\subsection{Continuation and adjoint}

\begingroup
\small
\setlength{\tabcolsep}{4pt}
\renewcommand{\arraystretch}{1.2}
\begin{longtable}{@{}p{0.78\textwidth} p{0.16\textwidth}@{}}
\toprule
\textbf{Citation} & \textbf{Status} \\
\midrule
\endhead
Keller 1977, pseudo-arclength continuation chapter & book chapter \\
Allgower \& Georg 1990/2003, \emph{Numerical Continuation Methods},
Springer / SIAM & textbook \\
Doedel et al.\ 2007, AUTO-07P manual & Concordia listing \\
Uecker 2021, \emph{Jahresber.~DMV} \textbf{123}, 199--248,
doi:10.1365/s13291-021-00241-5 & Verified \\
Hyon, Eisenberg, Liu 2010, \emph{Commun.~Math.~Sci.}\ \textbf{9}(2), 459--475 &
Verified \\
Schmuck, Bazant 2015, \emph{SIAM J.~Appl.~Math.}\ \textbf{75}(3), 1369--1401,
doi:10.1137/140968082 & Verified \\
Roy, Lin, Hahn 2024 EchemFEM, JOSS / LLNL-JRNL-860653 & Verified \\
Farrell, Ham, Funke, Rognes 2013, \emph{SIAM J.~Sci.~Comput.}\ \textbf{35}(4),
C369--C393, doi:10.1137/120873558 & Verified \\
Mitusch, Funke, Dokken 2019, \emph{J.~Open Source Softw.}\ \textbf{4}(38),
1292, doi:10.21105/joss.01292 & Verified \\
Plessix 2006, \emph{Geophys.~J.~Int.}\ \textbf{167}(2), 495--503,
doi:10.1111/j.1365-246X.2006.02978.x & Verified \\
Gunzburger 2003, \emph{Perspectives in Flow Control and Optimization}, SIAM &
textbook \\
Hinze, Pinnau, Ulbrich, Ulbrich 2009, \emph{Optimization with PDE Constraints},
Springer & textbook \\
Ascher \& Petzold 1998, \emph{Computer Methods for ODEs and DAEs}, SIAM &
textbook \\
\bottomrule
\end{longtable}
\endgroup

\subsection{Monolithic vs splitting PNP}

\begin{center}
\small
\begin{tabular}{@{}p{0.78\textwidth} p{0.16\textwidth}@{}}
\toprule
\textbf{Citation} & \textbf{Status} \\
\midrule
Gummel 1964, \emph{IEEE TED} \textbf{11}(10), 455--465 & original \\
Yu, Holst, McCammon 2007, \emph{J.~Comput.~Chem.}\ \textbf{28}, 1827--1839 &
monolithic PNP, biological \\
Liu \& Wang 2014, \emph{J.~Comput.~Phys.}\ \textbf{268}, 363--376 &
energy-stable FD \\
\bottomrule
\end{tabular}
\end{center}

\subsection{Not located / verification failed}

\begin{itemize}
\item \textbf{``Liu, Maxian, Shen, Zitelli, Waltz, Monteiro 2022, CMAME 396,
      115070, arXiv:2105.01163''} --- \textbf{fabricated}. arXiv:2105.01163
      resolves to Fu \& Xu 2022, CMAME 395, 115031. Replace.
\item \textbf{``Joshi, Seidel-Morgenstern, Tiller 2006''} --- not located in
      any database. If used in writeup, drop or replace with Brun--Reichert--K\"unsch
      2001 $+$ Quaiser--M\"onnigmann 2009 $+$ Pant 2018.
\end{itemize}

%% ============================================================
\section{Recommendations for the LaTeX Writeup's \texttt{thebibliography}}
%% ============================================================

\subsection{Critical corrections (blocking publication)}

\paragraph{Correction 1 --- Remove the fabricated CMAME 2022 reference;
replace with Fu \& Xu 2022.}

\begin{lstlisting}[style=bibstyle]
\bibitem{FuXu2022}
G.~Fu and Z.~Xu,
``High-order space-time finite element methods for the
Poisson--Nernst--Planck equations: positivity and unconditional energy
stability,''
{\em Computer Methods in Applied Mechanics and Engineering}, vol.~395,
art.~115031, 2022.
\href{https://doi.org/10.1016/j.cma.2022.115031}{doi:10.1016/j.cma.2022.115031}.
arXiv:2105.01163.
\end{lstlisting}

\paragraph{Correction 2 --- Disambiguate the Slotboom citation.}

\begin{lstlisting}[style=bibstyle]
\bibitem{Slotboom1969}
J.~W. Slotboom,
``Iterative scheme for 1- and 2-dimensional d.c.-transistor simulation,''
{\em Electronics Letters}, vol.~5, no.~26, pp.~677--678, 1969.
\href{https://doi.org/10.1049/el:19690510}{doi:10.1049/el:19690510}.
\end{lstlisting}

Add a one-sentence disambiguation in the methods text:
\begin{quote}
\small
``Slotboom variables and the log-density substitution are related by a
logarithm and additive shift --- $\ln(\Phi_n)=\ln(n)-\psi/\VT$ --- so they
share continuous structure but yield distinct Galerkin discretizations.''
\end{quote}

\paragraph{Correction 3 --- Disambiguate Zheng--Wei vs Zheng--Chen--Wei.}

\begin{lstlisting}[style=bibstyle]
\bibitem{ZhengWei2011}
Q.~Zheng and G.-W. Wei,
``Poisson--Boltzmann--Nernst--Planck model,''
{\em The Journal of Chemical Physics}, vol.~134, no.~19, art.~194101,
12~pp., 2011.
\href{https://doi.org/10.1063/1.3581031}{doi:10.1063/1.3581031}.
PMID:~21599038.

\bibitem{ZhengChenWei2011}
Q.~Zheng, D.~Chen, and G.-W. Wei,
``Second-order Poisson--Nernst--Planck solver for ion transport,''
{\em Journal of Computational Physics}, vol.~230, no.~13,
pp.~5239--5262, 2011.
\href{https://doi.org/10.1016/j.jcp.2011.03.020}{doi:10.1016/j.jcp.2011.03.020}.
\end{lstlisting}

Do \textbf{not} cite Zheng--Chen--Wei 2011 as the PBNP source.

\subsection{Recommended additions}

\begin{lstlisting}[style=bibstyle]
\bibitem{MettiXuLiu2016}
M.~S. Metti, J.~Xu, and C.~Liu,
``Energetically stable discretizations for charge transport and
electrokinetic models,''
{\em Journal of Computational Physics}, vol.~306, pp.~1--18, 2016.
\href{https://doi.org/10.1016/j.jcp.2015.10.053}{doi:10.1016/j.jcp.2015.10.053}.

\bibitem{CancesEtAl2022}
C.~Canc\`es, C.~Chainais-Hillairet, B.~Gaudeul, and J.~Fuhrmann,
``Entropy and convergence analysis for two finite-volume schemes for a
Nernst--Planck--Poisson system with ion-volume constraints,''
{\em Numerische Mathematik}, vol.~151, pp.~949--986, 2022.
\href{https://doi.org/10.1007/s00211-022-01279-y}{doi:10.1007/s00211-022-01279-y}.

\bibitem{Tafel1905}
J.~Tafel,
``\"Uber die Polarisation bei kathodischer Wasserstoffentwicklung,''
{\em Zeitschrift f\"ur Physikalische Chemie}, vol.~50U, no.~1,
pp.~641--712, 1905.
\href{https://doi.org/10.1515/zpch-1905-5043}{doi:10.1515/zpch-1905-5043}.

\bibitem{FattalKupferman2004}
R.~Fattal and R.~Kupferman,
``Constitutive laws for the matrix-logarithm of the conformation tensor,''
{\em Journal of Non-Newtonian Fluid Mechanics}, vol.~123, no.~2--3,
pp.~281--285, 2004.
\href{https://doi.org/10.1016/j.jnnfm.2004.08.008}{doi:10.1016/j.jnnfm.2004.08.008}.

\bibitem{LuHolstMcCammonZhou2010}
B.~Lu, M.~J. Holst, J.~A. McCammon, and Y.~C. Zhou,
``Poisson--Nernst--Planck equations for simulating biomolecular
diffusion-reaction processes I: Finite element solutions,''
{\em Journal of Computational Physics}, vol.~229, pp.~6979--6994, 2010.
\href{https://doi.org/10.1016/j.jcp.2010.05.035}{doi:10.1016/j.jcp.2010.05.035}.

\bibitem{Wroblowa1976}
H.~Wroblowa, Yen-Chi-Pan, and G.~Razumney,
``Electroreduction of oxygen: a new mechanistic criterion,''
{\em Journal of Electroanalytical Chemistry and Interfacial
Electrochemistry}, vol.~69, no.~2, pp.~195--201, 1976.
\href{https://doi.org/10.1016/S0022-0728(76)80250-1}{doi:10.1016/S0022-0728(76)80250-1}.

\bibitem{KulkarniSiahrostami2018}
A.~Kulkarni, S.~Siahrostami, A.~Patel, and J.~K. N{\o}rskov,
``Understanding catalytic activity trends in the oxygen reduction reaction,''
{\em Chemical Reviews}, vol.~118, no.~5, pp.~2302--2312, 2018.
\href{https://doi.org/10.1021/acs.chemrev.7b00488}{doi:10.1021/acs.chemrev.7b00488}.

\bibitem{FarrellHamFunkeRognes2013}
P.~E. Farrell, D.~A. Ham, S.~W. Funke, and M.~E. Rognes,
``Automated derivation of the adjoint of high-level transient finite
element programs,''
{\em SIAM Journal on Scientific Computing}, vol.~35, no.~4,
pp.~C369--C393, 2013.
\href{https://doi.org/10.1137/120873558}{doi:10.1137/120873558}.

\bibitem{MituschFunkeDokken2019}
S.~K. Mitusch, S.~W. Funke, and J.~S. Dokken,
``dolfin-adjoint 2018.1: automated adjoints for FEniCS and Firedrake,''
{\em Journal of Open Source Software}, vol.~4, no.~38, p.~1292, 2019.
\href{https://doi.org/10.21105/joss.01292}{doi:10.21105/joss.01292}.

\bibitem{MarkowichRinghoferSchmeiser1990}
P.~A. Markowich, C.~A. Ringhofer, and C.~Schmeiser,
{\em Semiconductor Equations}.
Wien: Springer, 1990.
\href{https://doi.org/10.1007/978-3-7091-6961-2}{doi:10.1007/978-3-7091-6961-2}.

\bibitem{KilicBazantAjdari2007a}
M.~S. Kilic, M.~Z. Bazant, and A.~Ajdari,
``Steric effects in the dynamics of electrolytes at large applied
voltages. I. Double-layer charging,''
{\em Physical Review E}, vol.~75, art.~021502, 2007.
\href{https://doi.org/10.1103/PhysRevE.75.021502}{doi:10.1103/PhysRevE.75.021502}.

\bibitem{BrunReichertKunsch2001}
R.~Brun, P.~Reichert, and H.~R. K\"unsch,
``Practical identifiability analysis of large environmental simulation
models,''
{\em Water Resources Research}, vol.~37, no.~4, pp.~1015--1030, 2001.
\href{https://doi.org/10.1029/2000WR900350}{doi:10.1029/2000WR900350}.

\bibitem{Pant2018}
S.~Pant,
``Information sensitivity functions to assess parameter information
gain and identifiability of dynamical systems,''
{\em Journal of the Royal Society Interface}, vol.~15,
art.~20170871, 2018.
\href{https://doi.org/10.1098/rsif.2017.0871}{doi:10.1098/rsif.2017.0871}.
\end{lstlisting}

\subsection{Recommended deletions}

\begin{itemize}
\item The fabricated \textbf{Liu--Maxian--Shen--Zitelli--Waltz--Monteiro 2022}
      entry must be removed.
\item If the writeup cites \textbf{Joshi, Seidel-Morgenstern, Tiller 2006},
      drop or replace with Brun~2001 $+$ Quaiser~2009 $+$ Pant~2018.
\item The misattributed Slotboom citation as printed (chimera 1973$+$1977
      conflation) must be replaced with Slotboom~1969.
\item Reassign Zheng--Chen--Wei 2011 JCP from PBNP-source role to
      full-PNP-numerical-methods role.
\end{itemize}

\subsection{Code-pointer corrections (writeup prose, not bibliography)}

\begin{enumerate}
\item \textbf{Line range for the log-rate branch}: change
      \texttt{forms\_logc.py:294-360} to \texttt{forms\_logc.py:294-388}.
\item \textbf{Cathodic-vs-anodic algebra description}: tighten ``both branches
      subtracted with $\sum_f\nu_f(u_{\mathrm{sp}(f)}-\ln c_{\mathrm{ref},f})$''
      to specify ``the cathodic branch carries the full stoichiometric-power
      loop (\texttt{forms\_logc.py:331-337}); the anodic branch is
      single-species (lines 343--345 and fallback 350--353). For the
      production setup $\mathrm{R}_2$ is irreversible
      (\texttt{v24\_3sp\_logc\_vs\_4sp\_validation.py:233}), so its anodic
      branch is zero, and the algebraic asymmetry has no production-output
      consequence.''
\item \textbf{Architectural note on \texttt{add\_boltzmann()}}: the writeup's
      pointer ``see \texttt{add\_boltzmann()} in the current study scripts''
      is technically correct but architecturally fragile (the function is
      monkey-patched in 7 study scripts). Tightening recommendation: future
      work should migrate the Boltzmann contribution into
      \texttt{Forward/bv\_solver/forms\_logc.build\_forms\_logc()} as an
      option.
\end{enumerate}

%% ============================================================
\section*{Document Provenance}
%% ============================================================
This document was synthesized from four parallel research investigations
(two web searches, one academic-literature search, one codebase audit) plus
direct DOI/arXiv verification of every cited paper. Raw research outputs are
preserved at \path{.research/pnp-bv-literature-provenance/}
(\texttt{web-logc.md}, \texttt{web-pbnp-lograte.md}, \texttt{lit-context.md},
\texttt{codebase.md}, \texttt{SUMMARY.md}).

\end{document}
